\scanpage{409}\mspage{148}
\deleted{\unsure{en d'autres}}

En d'autres termes, on peut \unsure{formellmt|formellement} regarder
la catégorie somme
\[
\underline{\mathrm{Ch}}_{\pm}^{\deleted{\unsure{W}}}(X)=
\underline{\mathrm{Ch}}_+^{\deleted{W}}(X)\amalg\underline{\mathrm{Ch}}_-^{\deleted{W}}(X)
\]
et \deleted{\unsure{T}}l'opération translation $T$ comme un
automorphisme qui échange les deux
morceaux, et $T^p$ n'est autre que son
itéré $p$-ième, qui échange les morceaux $+$ et $-$
si $p$ impair, qui les stabilise si $p$ pair.

Il y a aussi l'automorphisme \inserted{involutif} de \deleted{\unsure{symétrie}} \inserted{passage}
au chemin opposé, qui échange source et but,
et que je définis (conformément à l'isom
canonique identique
\[
\underline t\xrightarrow{\sim}\underline t^{\circ}\quad)
\]
comme la symétrie $S$
\[
\left\{\begin{aligned}
&\deleted{\unsure{\mathrm{opp}}}S\bigl((a_i)(u_i)\bigr)=\bigl((b_j),(\unsure{u_j|v_j})\bigr)\qquad\text{avec}\\
&b_j=a_{-j}\\
&u_j\deleted{\unsure{-}}=u_{-j+1}
\end{aligned}\right.
\]
Si je le regarde comme \unsure{opérant} sur $\underline{\mathrm{Ch}}_{\pm}^{W}(X)$,
il est lié \deleted{\unsure{à}}\inserted{a}u foncteur $T$ par
\[
\deleted{\unsure{TS}}\left\{\begin{aligned}S^2&=\mathrm{id}\\STS^{-1}&=T^{-1}\end{aligned}\right.,
\]
de sorte que c'est le groupe $\text{½}$ direct de
$\mathbb Z/2\unsure{\mathbb Z|0}$ par $\mathbb Z$ (groupe affine de dim 1 \unsure{sur}
l'anneau $\mathbb Z$) qui opère sur $\underline{\mathrm{Ch}}_{\pm}$.
% archive sheet number: 408
